
\documentclass{article}
\usepackage{colortbl}
\usepackage{makecell}
\usepackage{multirow}
\usepackage{supertabular}

\begin{document}

\newcounter{utterance}

\centering \large Interaction Transcript for game `clean\_up', experiment `0\_easy\_3obj\_et', episode 1 with Qwen3.6{-}35B{-}A3B{-}FP8{-}without{-}reasoning.
\vspace{24pt}

{ \footnotesize  \setcounter{utterance}{1}
\setlength{\tabcolsep}{0pt}
\begin{supertabular}{c@{$\;$}|p{.15\linewidth}@{}p{.15\linewidth}p{.15\linewidth}p{.15\linewidth}p{.15\linewidth}p{.15\linewidth}}
    \# & $\;$A & \multicolumn{4}{c}{Game Master} & $\;\:$B\\
    \hline

    \theutterance \stepcounter{utterance}  
    & & \multicolumn{4}{p{0.6\linewidth}}{
        \cellcolor[rgb]{0.9,0.9,0.9}{
            \makecell[{{p{\linewidth}}}]{
                \texttt{\tiny{[P1$\langle$GM]}}
                \texttt{Ma olen teie mängujuht ja te mängite ühismängu järgmise ruudustikuga mängulauana:} \\
\\ 
\texttt{```} \\
\texttt{ 1234567} \\
\texttt{╔═══════╗} \\
\texttt{║◌◌◌◌◌◌◌║ 1} \\
\texttt{║◌C◌◌◌◌◌║ 2} \\
\texttt{║◌◌◌◌◌┌─╢ 3} \\
\texttt{║◌◌◌┌─┤◌║ 4} \\
\texttt{║◌◌◌│L├─╢ 5} \\
\texttt{╟───┼─┘◌║ 6} \\
\texttt{║◌◌◌│◌◌P║ 7} \\
\texttt{╚═══╧═══╝} \\
\\ 
\texttt{```} \\
\\ 
\texttt{* Ülemisel serval on x{-}koordinaadid, mis suurenevad paremale, ja paremal serval y{-}koordinaadid, mis suurenevad allapoole.} \\
\texttt{* Järgmised objektid on teie ruudustikule juhuslikult paigutatud: 'C', 'L', 'P'.} \\
\\ 
\\ 
\texttt{Teine mängija näeb mängulaua varianti, kus objektid on paigutatud erinevatesse juhuslikesse asukohtadesse. Te ei näe teise mängija mängulauda ja tema ei näe teie mängulauda.} \\
\\ 
\\ 
\texttt{**Eesmärk:**} \\
\\ 
\texttt{Mõlemad mängijad peavad oma taustal olevaid objekte liigutama nii, et identsed objektid satuksid samadesse koordinaatidesse. Te peate teise mängijaga suhtlema, et kokku leppida ühises eesmärgi konfiguratsioonis.} \\
\\ 
\\ 
\texttt{**Reeglid:**} \\
\\ 
\texttt{* Igal käigul saate saata täpselt ühe kahest järgmisest käsust:} \\
\texttt{ 1. `ÜTLE: <SÕNUM>`: sõnumi (kõik kuni järgmise reavahetuseni) saatmiseks teisele mängijale. Ma edastan selle teie partnerile.} \\
\texttt{2. `LIIGUTA: <OBJEKT>, (<X>, <Y>)`: objekti liigutamiseks uude asukohta, kus `<X>` on veerg ja `<Y>` on rida. Ma teatan teile, kas teie käik õnnestus või mitte.} \\
\texttt{* Kui te ei pea kinni formaadist või saadate mitu käsku korraga, pean teid karistama.} \\
\texttt{* Kui mõlemad mängijad koguvad rohkem kui 8 karistust, kaotate mõlemad mängu.} \\
\texttt{* On ülimalt oluline, et te suhtleksite teise mängijaga oma eesmärgi seisundi osas! *Ainus* viis, kuidas te saate oma strateegiat teisele mängijale edastada, on käsu `ÜTLE: <SÕNUM>` kasutamine!} \\
\\ 
\\ 
\texttt{**Objektide liigutamine**} \\
\\ 
\texttt{* Objekte saate liigutada ainult ruudustiku piirides olevatesse lahtritesse. Sihtlahter peab olema tühi, st see peab sisaldama ainult sümbolit '◌'.} \\
\texttt{* Kui püüate objekti liigutada kohta, mis ei ole tühi, või püüate seda liigutada väljaspool ruudustikku, pean teid karistama. Saate uue katse.} \\
\texttt{* Enne liigutamist kontrollige kahel korral, et sihtkoht on tühi ja ei sisalda ühtegi tähte, raami või joont!} \\
\\ 
\\ 
\texttt{**Mängu lõpp**} \\
\\ 
\texttt{Kui arvate, et olete saavutanud kõigi objektide joondamise eesmärgi, võite paluda teisel mängijal mäng lõpetada, saates `ÜTLE: valmis?`. Kui teine mängija palub teil mängu lõpetada ja te vastate `ÜTLE: valmis!`, siis mäng lõpeb.} \\
\\ 
\texttt{Mõlemad mängijad võidavad, kui mäng lõpeb 12 vooru jooksul, kus üks voor on määratletud kahe mängija poolt ühe kehtiva käsu saatmisena.} \\
\\ 
\texttt{**Punktiarvestus:**} \\
\\ 
\texttt{Mida lähemal on identsed objektid mõlemal mängulaual, seda rohkem punkte saate. Karistused vähendavad teie punkte. Kas suudate rekordit ületada?} \\
\\ 
\\ 
\texttt{Palun saada teisele mängijale sõnum, et mängu alustada!} \\
            }
        }
    }
    & \\ \\

    \theutterance \stepcounter{utterance}  
    & \multicolumn{4}{p{0.6\linewidth}}{
        \cellcolor[rgb]{0.9,0.9,0.9}{
            \makecell[{{p{\linewidth}}}]{
                \texttt{\tiny{[P1$\rangle$GM]}}
                \texttt{ÜTLE: Tere! Alustame mängu. Minu laual on objektid järgmistes kohtades: C(2,2), L(4,5), P(6,7). Palun kirjuta oma objektide asukohad, et saaksime kokku leppida lõppkonfiguratsioonis.} \\
            }
        }
    }
    & & \\ \\

    \theutterance \stepcounter{utterance}  
    & & & \multicolumn{2}{p{0.3\linewidth}}{
        \cellcolor[rgb]{0.9,0.9,0.9}{
            \makecell[{{p{\linewidth}}}]{
                \texttt{\tiny{[GM$|$GM]}}
                \texttt{ÜTLE: Tere! Alustame mängu. Minu laual on objektid järgmistes kohtades: C(2,2), L(4,5), P(6,7). Palun kirjuta oma objektide asukohad, et saaksime kokku leppida lõppkonfiguratsioonis.} \\
            }
        }
    }
    & & \\ \\

    \theutterance \stepcounter{utterance}  
    & & \multicolumn{4}{p{0.6\linewidth}}{
        \cellcolor[rgb]{0.9,0.9,0.9}{
            \makecell[{{p{\linewidth}}}]{
                \texttt{\tiny{[GM$\rangle$P2]}}
                \texttt{Ma olen teie mängujuht ja te mängite ühismängu järgmise ruudustikuga mängulauana:} \\
\\ 
\texttt{```} \\
\texttt{ 1234567} \\
\texttt{╔═══════╗} \\
\texttt{║◌◌◌◌◌◌◌║ 1} \\
\texttt{║◌◌C◌◌◌◌║ 2} \\
\texttt{║◌◌◌◌◌┌─╢ 3} \\
\texttt{║◌◌◌┌─┤◌║ 4} \\
\texttt{║◌P◌│◌├─╢ 5} \\
\texttt{╟───┼─┘◌║ 6} \\
\texttt{║L◌◌│◌◌◌║ 7} \\
\texttt{╚═══╧═══╝} \\
\\ 
\texttt{```} \\
\\ 
\texttt{* Ülemisel serval on x{-}koordinaadid, mis suurenevad paremale, ja paremal serval y{-}koordinaadid, mis suurenevad allapoole.} \\
\texttt{* Järgmised objektid on teie ruudustikule juhuslikult paigutatud: 'C', 'L', 'P'.} \\
\\ 
\\ 
\texttt{Teine mängija näeb mängulaua varianti, kus objektid on paigutatud erinevatesse juhuslikesse asukohtadesse. Te ei näe teise mängija mängulauda ja tema ei näe teie mängulauda.} \\
\\ 
\\ 
\texttt{**Eesmärk:**} \\
\\ 
\texttt{Mõlemad mängijad peavad oma taustal olevaid objekte liigutama nii, et identsed objektid satuksid samadesse koordinaatidesse. Te peate teise mängijaga suhtlema, et kokku leppida ühises eesmärgi konfiguratsioonis.} \\
\\ 
\\ 
\texttt{**Reeglid:**} \\
\\ 
\texttt{* Igal käigul saate saata täpselt ühe kahest järgmisest käsust:} \\
\texttt{ 1. `ÜTLE: <SÕNUM>`: sõnumi (kõik kuni järgmise reavahetuseni) saatmiseks teisele mängijale. Ma edastan selle teie partnerile.} \\
\texttt{2. `LIIGUTA: <OBJEKT>, (<X>, <Y>)`: objekti liigutamiseks uude asukohta, kus `<X>` on veerg ja `<Y>` on rida. Ma teatan teile, kas teie käik õnnestus või mitte.} \\
\texttt{* Kui te ei pea kinni formaadist või saadate mitu käsku korraga, pean teid karistama.} \\
\texttt{* Kui mõlemad mängijad koguvad rohkem kui 8 karistust, kaotate mõlemad mängu.} \\
\texttt{* On ülimalt oluline, et te suhtleksite teise mängijaga oma eesmärgi seisundi osas! *Ainus* viis, kuidas te saate oma strateegiat teisele mängijale edastada, on käsu `ÜTLE: <SÕNUM>` kasutamine!} \\
\\ 
\\ 
\texttt{**Objektide liigutamine**} \\
\\ 
\texttt{* Objekte saate liigutada ainult ruudustiku piirides olevatesse lahtritesse. Sihtlahter peab olema tühi, st see peab sisaldama ainult sümbolit '◌'.} \\
\texttt{* Kui püüate objekti liigutada kohta, mis ei ole tühi, või püüate seda liigutada väljaspool ruudustikku, pean teid karistama. Saate uue katse.} \\
\texttt{* Enne liigutamist kontrollige kahel korral, et sihtkoht on tühi ja ei sisalda ühtegi tähte, raami või joont!} \\
\\ 
\\ 
\texttt{**Mängu lõpp**} \\
\\ 
\texttt{Kui arvate, et olete saavutanud kõigi objektide joondamise eesmärgi, võite paluda teisel mängijal mäng lõpetada, saates `ÜTLE: valmis?`. Kui teine mängija palub teil mängu lõpetada ja te vastate `ÜTLE: valmis!`, siis mäng lõpeb.} \\
\\ 
\texttt{Mõlemad mängijad võidavad, kui mäng lõpeb 12 vooru jooksul, kus üks voor on määratletud kahe mängija poolt ühe kehtiva käsu saatmisena.} \\
\\ 
\texttt{**Punktiarvestus:**} \\
\\ 
\texttt{Mida lähemal on identsed objektid mõlemal mängulaual, seda rohkem punkte saate. Karistused vähendavad teie punkte. Kas suudate rekordit ületada?} \\
\\ 
\texttt{Teine mängija alustas mängu selle sõnumi saatmisega:} \\
\\ 
\texttt{"Tere! Alustame mängu. Minu laual on objektid järgmistes kohtades: C(2,2), L(4,5), P(6,7). Palun kirjuta oma objektide asukohad, et saaksime kokku leppida lõppkonfiguratsioonis."} \\
\\ 
\texttt{Mis on sinu esimene käsk?} \\
            }
        }
    }
    & \\ \\

    \theutterance \stepcounter{utterance}  
    & & & \multicolumn{4}{p{0.6\linewidth}}{
        \cellcolor[rgb]{0.9,0.9,0.9}{
            \makecell[{{p{\linewidth}}}]{
                \texttt{\tiny{[GM$\langle$P2]}}
                \texttt{Tere! Minu laual on objektid järgmistes kohtades: C(3,2), L(1,7), P(2,5).} \\
            }
        }
    }
     \\ \\

    \theutterance \stepcounter{utterance}  
    & & & \multicolumn{2}{p{0.3\linewidth}}{
        \cellcolor[rgb]{0.9,0.9,0.9}{
            \makecell[{{p{\linewidth}}}]{
                \texttt{\tiny{[GM$|$GM]}}
                \texttt{Tere! Minu laual on objektid järgmistes kohtades: C(3,2), L(1,7), P(2,5).} \\
            }
        }
    }
    & & \\ \\

    \theutterance \stepcounter{utterance}  
    & & & \multicolumn{2}{p{0.3\linewidth}}{
        \cellcolor[rgb]{0.9,0.9,0.9}{
            \makecell[{{p{\linewidth}}}]{
                \texttt{\tiny{[GM$|$GM]}}
                \texttt{Invalid response format} \\
            }
        }
    }
    & & \\ \\

    \theutterance \stepcounter{utterance}  
    & & \multicolumn{4}{p{0.6\linewidth}}{
        \cellcolor[rgb]{0.9,0.9,0.9}{
            \makecell[{{p{\linewidth}}}]{
                \texttt{\tiny{[GM$\rangle$P2]}}
                \texttt{Karistus: Teie sõnum ei ole oodatud formaadis!} \\
\\ 
\texttt{Veendu, et sinu vastus sisaldab ainult kas `ÜTLE: <SÕNUM>` või `LIIGUTA: <OBJEKT>, (<X>, <Y>)` ja mitte midagi muud!} \\
\\ 
\texttt{Teil on kollektiivselt kogunenud 1 karistust 8 karistusest.} \\
\texttt{Palun proovi uuesti!} \\
            }
        }
    }
    & \\ \\

    \theutterance \stepcounter{utterance}  
    & & & \multicolumn{4}{p{0.6\linewidth}}{
        \cellcolor[rgb]{0.9,0.9,0.9}{
            \makecell[{{p{\linewidth}}}]{
                \texttt{\tiny{[GM$\langle$P2]}}
                \texttt{ÜTLE: Tere! Minu laual on objektid järgmistes kohtades: C(3,2), L(1,7), P(2,5).} \\
            }
        }
    }
     \\ \\

    \theutterance \stepcounter{utterance}  
    & & & \multicolumn{2}{p{0.3\linewidth}}{
        \cellcolor[rgb]{0.9,0.9,0.9}{
            \makecell[{{p{\linewidth}}}]{
                \texttt{\tiny{[GM$|$GM]}}
                \texttt{ÜTLE: Tere! Minu laual on objektid järgmistes kohtades: C(3,2), L(1,7), P(2,5).} \\
            }
        }
    }
    & & \\ \\

    \theutterance \stepcounter{utterance}  
    & & \multicolumn{4}{p{0.6\linewidth}}{
        \cellcolor[rgb]{0.9,0.9,0.9}{
            \makecell[{{p{\linewidth}}}]{
                \texttt{\tiny{[P1$\langle$GM]}}
                \texttt{Sinu sõnum on teisele mängijale edastatud.} \\
\\ 
\\ 
\texttt{Sa mängid praegu vooru 1 maksimaalsest 12 voorust.} \\
\texttt{Teil on kollektiivselt kogunenud 1 karistust 8 karistusest.} \\
\texttt{Teine mängija saatis selle sõnumi:} \\
\\ 
\\ 
\texttt{"Tere! Minu laual on objektid järgmistes kohtades: C(3,2), L(1,7), P(2,5)."} \\
\\ 
\\ 
\texttt{Mis on sinu järgmine käsk?} \\
            }
        }
    }
    & \\ \\

    \theutterance \stepcounter{utterance}  
    & \multicolumn{4}{p{0.6\linewidth}}{
        \cellcolor[rgb]{0.9,0.9,0.9}{
            \makecell[{{p{\linewidth}}}]{
                \texttt{\tiny{[P1$\rangle$GM]}}
                \texttt{ÜTLE: Suurepärane! Nüüd teame mõlema asukohti. Et lihtsustada liikumist, pakun välja järgmise eesmärgi: viime kõik objektid võimalikult lähedale ühte keskpunkti või lihtsalt kokku ühte nurka. Kuna minu C on (2,2) ja sinu C on (3,2), on need juba kõrvuti. Samamoodi on meie L ja P suhteliselt lähedal.} \\
\\ 
\texttt{Kas oled nõus, et proovime liigutada kõik objektid ritta või veergu? Näiteks võiksime püüda viia kõik objektid veergu 2 või ritta 2? Või on sul parem idee? Palun anna teada oma strateegiast, enne kui hakkame liikuma, et vältida karistusi.} \\
            }
        }
    }
    & & \\ \\

    \theutterance \stepcounter{utterance}  
    & & & \multicolumn{2}{p{0.3\linewidth}}{
        \cellcolor[rgb]{0.9,0.9,0.9}{
            \makecell[{{p{\linewidth}}}]{
                \texttt{\tiny{[GM$|$GM]}}
                \texttt{ÜTLE: Suurepärane! Nüüd teame mõlema asukohti. Et lihtsustada liikumist, pakun välja järgmise eesmärgi: viime kõik objektid võimalikult lähedale ühte keskpunkti või lihtsalt kokku ühte nurka. Kuna minu C on (2,2) ja sinu C on (3,2), on need juba kõrvuti. Samamoodi on meie L ja P suhteliselt lähedal.} \\
\\ 
\texttt{Kas oled nõus, et proovime liigutada kõik objektid ritta või veergu? Näiteks võiksime püüda viia kõik objektid veergu 2 või ritta 2? Või on sul parem idee? Palun anna teada oma strateegiast, enne kui hakkame liikuma, et vältida karistusi.} \\
            }
        }
    }
    & & \\ \\

    \theutterance \stepcounter{utterance}  
    & & & \multicolumn{2}{p{0.3\linewidth}}{
        \cellcolor[rgb]{0.9,0.9,0.9}{
            \makecell[{{p{\linewidth}}}]{
                \texttt{\tiny{[GM$|$GM]}}
                \texttt{Invalid format: tail is not empty: '} \\
\\ 
\texttt{Kas oled nõus, et proovime liigutada kõik objektid ritta või veergu? Näiteks võiksime püüda viia kõik objektid veergu 2 või ritta 2? Või on sul parem idee? Palun anna teada oma strateegiast, enne kui hakkame liikuma, et vältida karistusi.} \\
            }
        }
    }
    & & \\ \\

    \theutterance \stepcounter{utterance}  
    & & \multicolumn{4}{p{0.6\linewidth}}{
        \cellcolor[rgb]{0.9,0.9,0.9}{
            \makecell[{{p{\linewidth}}}]{
                \texttt{\tiny{[P1$\langle$GM]}}
                \texttt{Karistus: Teie sõnum ei tohi sisaldada midagi pärast käsku!} \\
\\ 
\texttt{Veendu, et sinu vastus sisaldab ainult kas `ÜTLE: <SÕNUM>` või `LIIGUTA: <OBJEKT>, (<X>, <Y>)` ja mitte midagi muud!} \\
\\ 
\texttt{Teil on kollektiivselt kogunenud 2 karistust 8 karistusest.} \\
\texttt{Palun proovi uuesti!} \\
            }
        }
    }
    & \\ \\

    \theutterance \stepcounter{utterance}  
    & \multicolumn{4}{p{0.6\linewidth}}{
        \cellcolor[rgb]{0.9,0.9,0.9}{
            \makecell[{{p{\linewidth}}}]{
                \texttt{\tiny{[P1$\rangle$GM]}}
                \texttt{ÜTLE: Suurepärane! Nüüd teame mõlema asukohti. Et lihtsustada liikumist, pakun välja järgmise eesmärgi: viime kõik objektid võimalikult lähedale ühte keskpunkti või lihtsalt kokku ühte nurka. Kuna minu C on (2,2) ja sinu C on (3,2), on need juba kõrvuti. Samamoodi on meie L ja P suhteliselt lähedal. Kas oled nõus, et proovime liigutada kõik objektid ritta või veergu? Näiteks võiksime püüda viia kõik objektid veergu 2 või ritta 2? Või on sul parem idee? Palun anna teada oma strateegiast, enne kui hakkame liikuma, et vältida karistusi.} \\
            }
        }
    }
    & & \\ \\

    \theutterance \stepcounter{utterance}  
    & & & \multicolumn{2}{p{0.3\linewidth}}{
        \cellcolor[rgb]{0.9,0.9,0.9}{
            \makecell[{{p{\linewidth}}}]{
                \texttt{\tiny{[GM$|$GM]}}
                \texttt{ÜTLE: Suurepärane! Nüüd teame mõlema asukohti. Et lihtsustada liikumist, pakun välja järgmise eesmärgi: viime kõik objektid võimalikult lähedale ühte keskpunkti või lihtsalt kokku ühte nurka. Kuna minu C on (2,2) ja sinu C on (3,2), on need juba kõrvuti. Samamoodi on meie L ja P suhteliselt lähedal. Kas oled nõus, et proovime liigutada kõik objektid ritta või veergu? Näiteks võiksime püüda viia kõik objektid veergu 2 või ritta 2? Või on sul parem idee? Palun anna teada oma strateegiast, enne kui hakkame liikuma, et vältida karistusi.} \\
            }
        }
    }
    & & \\ \\

    \theutterance \stepcounter{utterance}  
    & & \multicolumn{4}{p{0.6\linewidth}}{
        \cellcolor[rgb]{0.9,0.9,0.9}{
            \makecell[{{p{\linewidth}}}]{
                \texttt{\tiny{[GM$\rangle$P2]}}
                \texttt{Sinu sõnum on teisele mängijale edastatud.} \\
\\ 
\\ 
\texttt{Sa mängid praegu vooru 2 maksimaalsest 12 voorust.} \\
\texttt{Teil on kollektiivselt kogunenud 2 karistust 8 karistusest.} \\
\texttt{Teine mängija saatis selle sõnumi:} \\
\\ 
\\ 
\texttt{"Suurepärane! Nüüd teame mõlema asukohti. Et lihtsustada liikumist, pakun välja järgmise eesmärgi: viime kõik objektid võimalikult lähedale ühte keskpunkti või lihtsalt kokku ühte nurka. Kuna minu C on (2,2) ja sinu C on (3,2), on need juba kõrvuti. Samamoodi on meie L ja P suhteliselt lähedal. Kas oled nõus, et proovime liigutada kõik objektid ritta või veergu? Näiteks võiksime püüda viia kõik objektid veergu 2 või ritta 2? Või on sul parem idee? Palun anna teada oma strateegiast, enne kui hakkame liikuma, et vältida karistusi."} \\
\\ 
\\ 
\texttt{Mis on sinu järgmine käsk?} \\
            }
        }
    }
    & \\ \\

    \theutterance \stepcounter{utterance}  
    & & & \multicolumn{4}{p{0.6\linewidth}}{
        \cellcolor[rgb]{0.9,0.9,0.9}{
            \makecell[{{p{\linewidth}}}]{
                \texttt{\tiny{[GM$\langle$P2]}}
                \texttt{ÜTLE: Nõus, liigutame kõik objektid veergu 2, ritta 2. See on lihtne ja selge eesmärk. Alustame liikumisega.} \\
            }
        }
    }
     \\ \\

    \theutterance \stepcounter{utterance}  
    & & & \multicolumn{2}{p{0.3\linewidth}}{
        \cellcolor[rgb]{0.9,0.9,0.9}{
            \makecell[{{p{\linewidth}}}]{
                \texttt{\tiny{[GM$|$GM]}}
                \texttt{ÜTLE: Nõus, liigutame kõik objektid veergu 2, ritta 2. See on lihtne ja selge eesmärk. Alustame liikumisega.} \\
            }
        }
    }
    & & \\ \\

    \theutterance \stepcounter{utterance}  
    & & \multicolumn{4}{p{0.6\linewidth}}{
        \cellcolor[rgb]{0.9,0.9,0.9}{
            \makecell[{{p{\linewidth}}}]{
                \texttt{\tiny{[P1$\langle$GM]}}
                \texttt{Sinu sõnum on teisele mängijale edastatud.} \\
\\ 
\\ 
\texttt{Sa mängid praegu vooru 2 maksimaalsest 12 voorust.} \\
\texttt{Teil on kollektiivselt kogunenud 2 karistust 8 karistusest.} \\
\texttt{Teine mängija saatis selle sõnumi:} \\
\\ 
\\ 
\texttt{"Nõus, liigutame kõik objektid veergu 2, ritta 2. See on lihtne ja selge eesmärk. Alustame liikumisega."} \\
\\ 
\\ 
\texttt{Mis on sinu järgmine käsk?} \\
            }
        }
    }
    & \\ \\

    \theutterance \stepcounter{utterance}  
    & \multicolumn{4}{p{0.6\linewidth}}{
        \cellcolor[rgb]{0.9,0.9,0.9}{
            \makecell[{{p{\linewidth}}}]{
                \texttt{\tiny{[P1$\rangle$GM]}}
                \texttt{LIIGUTA: C, (2, 2)} \\
            }
        }
    }
    & & \\ \\

    \theutterance \stepcounter{utterance}  
    & & & \multicolumn{2}{p{0.3\linewidth}}{
        \cellcolor[rgb]{0.9,0.9,0.9}{
            \makecell[{{p{\linewidth}}}]{
                \texttt{\tiny{[GM$|$GM]}}
                \texttt{LIIGUTA: C, (2, 2)} \\
            }
        }
    }
    & & \\ \\

    \theutterance \stepcounter{utterance}  
    & & & \multicolumn{2}{p{0.3\linewidth}}{
        \cellcolor[rgb]{0.9,0.9,0.9}{
            \makecell[{{p{\linewidth}}}]{
                \texttt{\tiny{[GM$|$GM]}}
                \texttt{Karistus: (2,2) ei ole tühi, vaid sisaldab objekti 'C'.} \\
\\ 
\texttt{Teil on kollektiivselt kogunenud 3 karistust 8 karistusest.} \\
\texttt{Palun proovi uuesti!} \\
            }
        }
    }
    & & \\ \\

    \theutterance \stepcounter{utterance}  
    & & \multicolumn{4}{p{0.6\linewidth}}{
        \cellcolor[rgb]{0.9,0.9,0.9}{
            \makecell[{{p{\linewidth}}}]{
                \texttt{\tiny{[P1$\langle$GM]}}
                \texttt{Karistus: (2,2) ei ole tühi, vaid sisaldab objekti 'C'.} \\
\\ 
\texttt{Teil on kollektiivselt kogunenud 3 karistust 8 karistusest.} \\
\texttt{Palun proovi uuesti!} \\
            }
        }
    }
    & \\ \\

    \theutterance \stepcounter{utterance}  
    & \multicolumn{4}{p{0.6\linewidth}}{
        \cellcolor[rgb]{0.9,0.9,0.9}{
            \makecell[{{p{\linewidth}}}]{
                \texttt{\tiny{[P1$\rangle$GM]}}
                \texttt{LIIGUTA: C, (2, 1)} \\
            }
        }
    }
    & & \\ \\

    \theutterance \stepcounter{utterance}  
    & & & \multicolumn{2}{p{0.3\linewidth}}{
        \cellcolor[rgb]{0.9,0.9,0.9}{
            \makecell[{{p{\linewidth}}}]{
                \texttt{\tiny{[GM$|$GM]}}
                \texttt{LIIGUTA: C, (2, 1)} \\
            }
        }
    }
    & & \\ \\

    \theutterance \stepcounter{utterance}  
    & & & \multicolumn{2}{p{0.3\linewidth}}{
        \cellcolor[rgb]{0.9,0.9,0.9}{
            \makecell[{{p{\linewidth}}}]{
                \texttt{\tiny{[GM$|$GM]}}
                \texttt{C moved to (2, 1)} \\
            }
        }
    }
    & & \\ \\

    \theutterance \stepcounter{utterance}  
    & & \multicolumn{4}{p{0.6\linewidth}}{
        \cellcolor[rgb]{0.9,0.9,0.9}{
            \makecell[{{p{\linewidth}}}]{
                \texttt{\tiny{[GM$\rangle$P2]}}
                \texttt{Sinu sõnum on teisele mängijale edastatud.} \\
\\ 
\\ 
\texttt{Sa mängid praegu vooru 3 maksimaalsest 12 voorust.} \\
\texttt{Teil on kollektiivselt kogunenud 3 karistust 8 karistusest.} \\
\texttt{Teine mängija liigutas objekti oma ruudustikul.} \\
\\ 
\\ 
\texttt{Mis on sinu järgmine käsk?} \\
            }
        }
    }
    & \\ \\

    \theutterance \stepcounter{utterance}  
    & & & \multicolumn{4}{p{0.6\linewidth}}{
        \cellcolor[rgb]{0.9,0.9,0.9}{
            \makecell[{{p{\linewidth}}}]{
                \texttt{\tiny{[GM$\langle$P2]}}
                \texttt{ÜTLE: Liigutan oma objektid sihtkohta C(2,2), L(2,2), P(2,2). Kuna sihtkohad on samad, liigutan need järjest. Alustan C{-}ga.} \\
            }
        }
    }
     \\ \\

    \theutterance \stepcounter{utterance}  
    & & & \multicolumn{2}{p{0.3\linewidth}}{
        \cellcolor[rgb]{0.9,0.9,0.9}{
            \makecell[{{p{\linewidth}}}]{
                \texttt{\tiny{[GM$|$GM]}}
                \texttt{ÜTLE: Liigutan oma objektid sihtkohta C(2,2), L(2,2), P(2,2). Kuna sihtkohad on samad, liigutan need järjest. Alustan C{-}ga.} \\
            }
        }
    }
    & & \\ \\

    \theutterance \stepcounter{utterance}  
    & & \multicolumn{4}{p{0.6\linewidth}}{
        \cellcolor[rgb]{0.9,0.9,0.9}{
            \makecell[{{p{\linewidth}}}]{
                \texttt{\tiny{[P1$\langle$GM]}}
                \texttt{Edukalt teisaldatud 'C' asukohta (2,1). Teie uuendatud ruudustik näeb välja selline:} \\
\texttt{```} \\
\texttt{ 1234567} \\
\texttt{╔═══════╗} \\
\texttt{║◌C◌◌◌◌◌║ 1} \\
\texttt{║◌◌◌◌◌◌◌║ 2} \\
\texttt{║◌◌◌◌◌┌─╢ 3} \\
\texttt{║◌◌◌┌─┤◌║ 4} \\
\texttt{║◌◌◌│L├─╢ 5} \\
\texttt{╟───┼─┘◌║ 6} \\
\texttt{║◌◌◌│◌◌P║ 7} \\
\texttt{╚═══╧═══╝} \\
\texttt{```} \\
\texttt{Sa mängid praegu vooru 3 maksimaalsest 12 voorust.} \\
\texttt{Teil on kollektiivselt kogunenud 3 karistust 8 karistusest.} \\
\texttt{Teine mängija saatis selle sõnumi:} \\
\\ 
\\ 
\texttt{"Liigutan oma objektid sihtkohta C(2,2), L(2,2), P(2,2). Kuna sihtkohad on samad, liigutan need järjest. Alustan C{-}ga."} \\
\\ 
\\ 
\texttt{Mis on sinu järgmine käsk?} \\
            }
        }
    }
    & \\ \\

    \theutterance \stepcounter{utterance}  
    & \multicolumn{4}{p{0.6\linewidth}}{
        \cellcolor[rgb]{0.9,0.9,0.9}{
            \makecell[{{p{\linewidth}}}]{
                \texttt{\tiny{[P1$\rangle$GM]}}
                \texttt{LIIGUTA: L, (2, 5)} \\
            }
        }
    }
    & & \\ \\

    \theutterance \stepcounter{utterance}  
    & & & \multicolumn{2}{p{0.3\linewidth}}{
        \cellcolor[rgb]{0.9,0.9,0.9}{
            \makecell[{{p{\linewidth}}}]{
                \texttt{\tiny{[GM$|$GM]}}
                \texttt{LIIGUTA: L, (2, 5)} \\
            }
        }
    }
    & & \\ \\

    \theutterance \stepcounter{utterance}  
    & & & \multicolumn{2}{p{0.3\linewidth}}{
        \cellcolor[rgb]{0.9,0.9,0.9}{
            \makecell[{{p{\linewidth}}}]{
                \texttt{\tiny{[GM$|$GM]}}
                \texttt{L moved to (2, 5)} \\
            }
        }
    }
    & & \\ \\

    \theutterance \stepcounter{utterance}  
    & & \multicolumn{4}{p{0.6\linewidth}}{
        \cellcolor[rgb]{0.9,0.9,0.9}{
            \makecell[{{p{\linewidth}}}]{
                \texttt{\tiny{[GM$\rangle$P2]}}
                \texttt{Sinu sõnum on teisele mängijale edastatud.} \\
\\ 
\\ 
\texttt{Sa mängid praegu vooru 4 maksimaalsest 12 voorust.} \\
\texttt{Teil on kollektiivselt kogunenud 3 karistust 8 karistusest.} \\
\texttt{Teine mängija liigutas objekti oma ruudustikul.} \\
\\ 
\\ 
\texttt{Mis on sinu järgmine käsk?} \\
            }
        }
    }
    & \\ \\

    \theutterance \stepcounter{utterance}  
    & & & \multicolumn{4}{p{0.6\linewidth}}{
        \cellcolor[rgb]{0.9,0.9,0.9}{
            \makecell[{{p{\linewidth}}}]{
                \texttt{\tiny{[GM$\langle$P2]}}
                \texttt{LIIGUTA: C, (2, 2)} \\
            }
        }
    }
     \\ \\

    \theutterance \stepcounter{utterance}  
    & & & \multicolumn{2}{p{0.3\linewidth}}{
        \cellcolor[rgb]{0.9,0.9,0.9}{
            \makecell[{{p{\linewidth}}}]{
                \texttt{\tiny{[GM$|$GM]}}
                \texttt{LIIGUTA: C, (2, 2)} \\
            }
        }
    }
    & & \\ \\

    \theutterance \stepcounter{utterance}  
    & & & \multicolumn{2}{p{0.3\linewidth}}{
        \cellcolor[rgb]{0.9,0.9,0.9}{
            \makecell[{{p{\linewidth}}}]{
                \texttt{\tiny{[GM$|$GM]}}
                \texttt{C moved to (2, 2)} \\
            }
        }
    }
    & & \\ \\

    \theutterance \stepcounter{utterance}  
    & & \multicolumn{4}{p{0.6\linewidth}}{
        \cellcolor[rgb]{0.9,0.9,0.9}{
            \makecell[{{p{\linewidth}}}]{
                \texttt{\tiny{[P1$\langle$GM]}}
                \texttt{Edukalt teisaldatud 'L' asukohta (2,5). Teie uuendatud ruudustik näeb välja selline:} \\
\texttt{```} \\
\texttt{ 1234567} \\
\texttt{╔═══════╗} \\
\texttt{║◌C◌◌◌◌◌║ 1} \\
\texttt{║◌◌◌◌◌◌◌║ 2} \\
\texttt{║◌◌◌◌◌┌─╢ 3} \\
\texttt{║◌◌◌┌─┤◌║ 4} \\
\texttt{║◌L◌│◌├─╢ 5} \\
\texttt{╟───┼─┘◌║ 6} \\
\texttt{║◌◌◌│◌◌P║ 7} \\
\texttt{╚═══╧═══╝} \\
\texttt{```} \\
\texttt{Sa mängid praegu vooru 4 maksimaalsest 12 voorust.} \\
\texttt{Teil on kollektiivselt kogunenud 3 karistust 8 karistusest.} \\
\texttt{Teine mängija liigutas objekti oma ruudustikul.} \\
\\ 
\\ 
\texttt{Mis on sinu järgmine käsk?} \\
            }
        }
    }
    & \\ \\

    \theutterance \stepcounter{utterance}  
    & \multicolumn{4}{p{0.6\linewidth}}{
        \cellcolor[rgb]{0.9,0.9,0.9}{
            \makecell[{{p{\linewidth}}}]{
                \texttt{\tiny{[P1$\rangle$GM]}}
                \texttt{LIIGUTA: P, (2, 7)} \\
            }
        }
    }
    & & \\ \\

    \theutterance \stepcounter{utterance}  
    & & & \multicolumn{2}{p{0.3\linewidth}}{
        \cellcolor[rgb]{0.9,0.9,0.9}{
            \makecell[{{p{\linewidth}}}]{
                \texttt{\tiny{[GM$|$GM]}}
                \texttt{LIIGUTA: P, (2, 7)} \\
            }
        }
    }
    & & \\ \\

    \theutterance \stepcounter{utterance}  
    & & & \multicolumn{2}{p{0.3\linewidth}}{
        \cellcolor[rgb]{0.9,0.9,0.9}{
            \makecell[{{p{\linewidth}}}]{
                \texttt{\tiny{[GM$|$GM]}}
                \texttt{P moved to (2, 7)} \\
            }
        }
    }
    & & \\ \\

    \theutterance \stepcounter{utterance}  
    & & \multicolumn{4}{p{0.6\linewidth}}{
        \cellcolor[rgb]{0.9,0.9,0.9}{
            \makecell[{{p{\linewidth}}}]{
                \texttt{\tiny{[GM$\rangle$P2]}}
                \texttt{Edukalt teisaldatud 'C' asukohta (2,2). Teie uuendatud ruudustik näeb välja selline:} \\
\texttt{```} \\
\texttt{ 1234567} \\
\texttt{╔═══════╗} \\
\texttt{║◌◌◌◌◌◌◌║ 1} \\
\texttt{║◌C◌◌◌◌◌║ 2} \\
\texttt{║◌◌◌◌◌┌─╢ 3} \\
\texttt{║◌◌◌┌─┤◌║ 4} \\
\texttt{║◌P◌│◌├─╢ 5} \\
\texttt{╟───┼─┘◌║ 6} \\
\texttt{║L◌◌│◌◌◌║ 7} \\
\texttt{╚═══╧═══╝} \\
\texttt{```} \\
\texttt{Sa mängid praegu vooru 5 maksimaalsest 12 voorust.} \\
\texttt{Teil on kollektiivselt kogunenud 3 karistust 8 karistusest.} \\
\texttt{Teine mängija liigutas objekti oma ruudustikul.} \\
\\ 
\\ 
\texttt{Mis on sinu järgmine käsk?} \\
            }
        }
    }
    & \\ \\

    \theutterance \stepcounter{utterance}  
    & & & \multicolumn{4}{p{0.6\linewidth}}{
        \cellcolor[rgb]{0.9,0.9,0.9}{
            \makecell[{{p{\linewidth}}}]{
                \texttt{\tiny{[GM$\langle$P2]}}
                \texttt{LIIGUTA: P, (2, 5)} \\
            }
        }
    }
     \\ \\

    \theutterance \stepcounter{utterance}  
    & & & \multicolumn{2}{p{0.3\linewidth}}{
        \cellcolor[rgb]{0.9,0.9,0.9}{
            \makecell[{{p{\linewidth}}}]{
                \texttt{\tiny{[GM$|$GM]}}
                \texttt{LIIGUTA: P, (2, 5)} \\
            }
        }
    }
    & & \\ \\

    \theutterance \stepcounter{utterance}  
    & & & \multicolumn{2}{p{0.3\linewidth}}{
        \cellcolor[rgb]{0.9,0.9,0.9}{
            \makecell[{{p{\linewidth}}}]{
                \texttt{\tiny{[GM$|$GM]}}
                \texttt{Karistus: (2,5) ei ole tühi, vaid sisaldab objekti 'P'.} \\
\\ 
\texttt{Teil on kollektiivselt kogunenud 4 karistust 8 karistusest.} \\
\texttt{Palun proovi uuesti!} \\
            }
        }
    }
    & & \\ \\

    \theutterance \stepcounter{utterance}  
    & & \multicolumn{4}{p{0.6\linewidth}}{
        \cellcolor[rgb]{0.9,0.9,0.9}{
            \makecell[{{p{\linewidth}}}]{
                \texttt{\tiny{[GM$\rangle$P2]}}
                \texttt{Karistus: (2,5) ei ole tühi, vaid sisaldab objekti 'P'.} \\
\\ 
\texttt{Teil on kollektiivselt kogunenud 4 karistust 8 karistusest.} \\
\texttt{Palun proovi uuesti!} \\
            }
        }
    }
    & \\ \\

    \theutterance \stepcounter{utterance}  
    & & & \multicolumn{4}{p{0.6\linewidth}}{
        \cellcolor[rgb]{0.9,0.9,0.9}{
            \makecell[{{p{\linewidth}}}]{
                \texttt{\tiny{[GM$\langle$P2]}}
                \texttt{LIIGUTA: P, (2, 6)} \\
            }
        }
    }
     \\ \\

    \theutterance \stepcounter{utterance}  
    & & & \multicolumn{2}{p{0.3\linewidth}}{
        \cellcolor[rgb]{0.9,0.9,0.9}{
            \makecell[{{p{\linewidth}}}]{
                \texttt{\tiny{[GM$|$GM]}}
                \texttt{LIIGUTA: P, (2, 6)} \\
            }
        }
    }
    & & \\ \\

    \theutterance \stepcounter{utterance}  
    & & & \multicolumn{2}{p{0.3\linewidth}}{
        \cellcolor[rgb]{0.9,0.9,0.9}{
            \makecell[{{p{\linewidth}}}]{
                \texttt{\tiny{[GM$|$GM]}}
                \texttt{Karistus: (2,6) ei ole tühi, vaid sisaldab objekti '─'.} \\
\\ 
\texttt{Teil on kollektiivselt kogunenud 5 karistust 8 karistusest.} \\
\texttt{Palun proovi uuesti!} \\
            }
        }
    }
    & & \\ \\

    \theutterance \stepcounter{utterance}  
    & & \multicolumn{4}{p{0.6\linewidth}}{
        \cellcolor[rgb]{0.9,0.9,0.9}{
            \makecell[{{p{\linewidth}}}]{
                \texttt{\tiny{[GM$\rangle$P2]}}
                \texttt{Karistus: (2,6) ei ole tühi, vaid sisaldab objekti '─'.} \\
\\ 
\texttt{Teil on kollektiivselt kogunenud 5 karistust 8 karistusest.} \\
\texttt{Palun proovi uuesti!} \\
            }
        }
    }
    & \\ \\

    \theutterance \stepcounter{utterance}  
    & & & \multicolumn{4}{p{0.6\linewidth}}{
        \cellcolor[rgb]{0.9,0.9,0.9}{
            \makecell[{{p{\linewidth}}}]{
                \texttt{\tiny{[GM$\langle$P2]}}
                \texttt{LIIGUTA: P, (2, 4)} \\
            }
        }
    }
     \\ \\

    \theutterance \stepcounter{utterance}  
    & & & \multicolumn{2}{p{0.3\linewidth}}{
        \cellcolor[rgb]{0.9,0.9,0.9}{
            \makecell[{{p{\linewidth}}}]{
                \texttt{\tiny{[GM$|$GM]}}
                \texttt{LIIGUTA: P, (2, 4)} \\
            }
        }
    }
    & & \\ \\

    \theutterance \stepcounter{utterance}  
    & & & \multicolumn{2}{p{0.3\linewidth}}{
        \cellcolor[rgb]{0.9,0.9,0.9}{
            \makecell[{{p{\linewidth}}}]{
                \texttt{\tiny{[GM$|$GM]}}
                \texttt{P moved to (2, 4)} \\
            }
        }
    }
    & & \\ \\

    \theutterance \stepcounter{utterance}  
    & & \multicolumn{4}{p{0.6\linewidth}}{
        \cellcolor[rgb]{0.9,0.9,0.9}{
            \makecell[{{p{\linewidth}}}]{
                \texttt{\tiny{[P1$\langle$GM]}}
                \texttt{Edukalt teisaldatud 'P' asukohta (2,7). Teie uuendatud ruudustik näeb välja selline:} \\
\texttt{```} \\
\texttt{ 1234567} \\
\texttt{╔═══════╗} \\
\texttt{║◌C◌◌◌◌◌║ 1} \\
\texttt{║◌◌◌◌◌◌◌║ 2} \\
\texttt{║◌◌◌◌◌┌─╢ 3} \\
\texttt{║◌◌◌┌─┤◌║ 4} \\
\texttt{║◌L◌│◌├─╢ 5} \\
\texttt{╟───┼─┘◌║ 6} \\
\texttt{║◌P◌│◌◌◌║ 7} \\
\texttt{╚═══╧═══╝} \\
\texttt{```} \\
\texttt{Sa mängid praegu vooru 5 maksimaalsest 12 voorust.} \\
\texttt{Teil on kollektiivselt kogunenud 5 karistust 8 karistusest.} \\
\texttt{Teine mängija liigutas objekti oma ruudustikul.} \\
\\ 
\\ 
\texttt{Mis on sinu järgmine käsk?} \\
            }
        }
    }
    & \\ \\

    \theutterance \stepcounter{utterance}  
    & \multicolumn{4}{p{0.6\linewidth}}{
        \cellcolor[rgb]{0.9,0.9,0.9}{
            \makecell[{{p{\linewidth}}}]{
                \texttt{\tiny{[P1$\rangle$GM]}}
                \texttt{ÜTLE: Meie objektid on nüüd veerus 2: C(2,1), L(2,5), P(2,7). Teine mängija ütleb, et viib oma objektid samasse kohta (2,2). See tähendab, et tema C, L ja P peaksid kõik olema lõpuks kohal (2,2). Meie objektid on juba veerus 2, kuid erinevatel ridadel. Et saavutada täielik kattumine, peaksime ka meie viima objektid (2,2){-}sse. Kuid (2,2) on praegu tühi. Liigutan C(2,1) kohale (2,2).} \\
            }
        }
    }
    & & \\ \\

    \theutterance \stepcounter{utterance}  
    & & & \multicolumn{2}{p{0.3\linewidth}}{
        \cellcolor[rgb]{0.9,0.9,0.9}{
            \makecell[{{p{\linewidth}}}]{
                \texttt{\tiny{[GM$|$GM]}}
                \texttt{ÜTLE: Meie objektid on nüüd veerus 2: C(2,1), L(2,5), P(2,7). Teine mängija ütleb, et viib oma objektid samasse kohta (2,2). See tähendab, et tema C, L ja P peaksid kõik olema lõpuks kohal (2,2). Meie objektid on juba veerus 2, kuid erinevatel ridadel. Et saavutada täielik kattumine, peaksime ka meie viima objektid (2,2){-}sse. Kuid (2,2) on praegu tühi. Liigutan C(2,1) kohale (2,2).} \\
            }
        }
    }
    & & \\ \\

    \theutterance \stepcounter{utterance}  
    & & \multicolumn{4}{p{0.6\linewidth}}{
        \cellcolor[rgb]{0.9,0.9,0.9}{
            \makecell[{{p{\linewidth}}}]{
                \texttt{\tiny{[GM$\rangle$P2]}}
                \texttt{Edukalt teisaldatud 'P' asukohta (2,4). Teie uuendatud ruudustik näeb välja selline:} \\
\texttt{```} \\
\texttt{ 1234567} \\
\texttt{╔═══════╗} \\
\texttt{║◌◌◌◌◌◌◌║ 1} \\
\texttt{║◌C◌◌◌◌◌║ 2} \\
\texttt{║◌◌◌◌◌┌─╢ 3} \\
\texttt{║◌P◌┌─┤◌║ 4} \\
\texttt{║◌◌◌│◌├─╢ 5} \\
\texttt{╟───┼─┘◌║ 6} \\
\texttt{║L◌◌│◌◌◌║ 7} \\
\texttt{╚═══╧═══╝} \\
\texttt{```} \\
\texttt{Sa mängid praegu vooru 6 maksimaalsest 12 voorust.} \\
\texttt{Teil on kollektiivselt kogunenud 5 karistust 8 karistusest.} \\
\texttt{Teine mängija saatis selle sõnumi:} \\
\\ 
\\ 
\texttt{"Meie objektid on nüüd veerus 2: C(2,1), L(2,5), P(2,7). Teine mängija ütleb, et viib oma objektid samasse kohta (2,2). See tähendab, et tema C, L ja P peaksid kõik olema lõpuks kohal (2,2). Meie objektid on juba veerus 2, kuid erinevatel ridadel. Et saavutada täielik kattumine, peaksime ka meie viima objektid (2,2){-}sse. Kuid (2,2) on praegu tühi. Liigutan C(2,1) kohale (2,2)."} \\
\\ 
\\ 
\texttt{Mis on sinu järgmine käsk?} \\
            }
        }
    }
    & \\ \\

    \theutterance \stepcounter{utterance}  
    & & & \multicolumn{4}{p{0.6\linewidth}}{
        \cellcolor[rgb]{0.9,0.9,0.9}{
            \makecell[{{p{\linewidth}}}]{
                \texttt{\tiny{[GM$\langle$P2]}}
                \texttt{LIIGUTA: P, (2, 2)} \\
            }
        }
    }
     \\ \\

    \theutterance \stepcounter{utterance}  
    & & & \multicolumn{2}{p{0.3\linewidth}}{
        \cellcolor[rgb]{0.9,0.9,0.9}{
            \makecell[{{p{\linewidth}}}]{
                \texttt{\tiny{[GM$|$GM]}}
                \texttt{LIIGUTA: P, (2, 2)} \\
            }
        }
    }
    & & \\ \\

    \theutterance \stepcounter{utterance}  
    & & & \multicolumn{2}{p{0.3\linewidth}}{
        \cellcolor[rgb]{0.9,0.9,0.9}{
            \makecell[{{p{\linewidth}}}]{
                \texttt{\tiny{[GM$|$GM]}}
                \texttt{Karistus: (2,2) ei ole tühi, vaid sisaldab objekti 'C'.} \\
\\ 
\texttt{Teil on kollektiivselt kogunenud 6 karistust 8 karistusest.} \\
\texttt{Palun proovi uuesti!} \\
            }
        }
    }
    & & \\ \\

    \theutterance \stepcounter{utterance}  
    & & \multicolumn{4}{p{0.6\linewidth}}{
        \cellcolor[rgb]{0.9,0.9,0.9}{
            \makecell[{{p{\linewidth}}}]{
                \texttt{\tiny{[GM$\rangle$P2]}}
                \texttt{Karistus: (2,2) ei ole tühi, vaid sisaldab objekti 'C'.} \\
\\ 
\texttt{Teil on kollektiivselt kogunenud 6 karistust 8 karistusest.} \\
\texttt{Palun proovi uuesti!} \\
            }
        }
    }
    & \\ \\

    \theutterance \stepcounter{utterance}  
    & & & \multicolumn{4}{p{0.6\linewidth}}{
        \cellcolor[rgb]{0.9,0.9,0.9}{
            \makecell[{{p{\linewidth}}}]{
                \texttt{\tiny{[GM$\langle$P2]}}
                \texttt{LIIGUTA: L, (2, 2)} \\
            }
        }
    }
     \\ \\

    \theutterance \stepcounter{utterance}  
    & & & \multicolumn{2}{p{0.3\linewidth}}{
        \cellcolor[rgb]{0.9,0.9,0.9}{
            \makecell[{{p{\linewidth}}}]{
                \texttt{\tiny{[GM$|$GM]}}
                \texttt{LIIGUTA: L, (2, 2)} \\
            }
        }
    }
    & & \\ \\

    \theutterance \stepcounter{utterance}  
    & & & \multicolumn{2}{p{0.3\linewidth}}{
        \cellcolor[rgb]{0.9,0.9,0.9}{
            \makecell[{{p{\linewidth}}}]{
                \texttt{\tiny{[GM$|$GM]}}
                \texttt{Karistus: (2,2) ei ole tühi, vaid sisaldab objekti 'C'.} \\
\\ 
\texttt{Teil on kollektiivselt kogunenud 7 karistust 8 karistusest.} \\
\texttt{Palun proovi uuesti!} \\
            }
        }
    }
    & & \\ \\

    \theutterance \stepcounter{utterance}  
    & & \multicolumn{4}{p{0.6\linewidth}}{
        \cellcolor[rgb]{0.9,0.9,0.9}{
            \makecell[{{p{\linewidth}}}]{
                \texttt{\tiny{[GM$\rangle$P2]}}
                \texttt{Karistus: (2,2) ei ole tühi, vaid sisaldab objekti 'C'.} \\
\\ 
\texttt{Teil on kollektiivselt kogunenud 7 karistust 8 karistusest.} \\
\texttt{Palun proovi uuesti!} \\
            }
        }
    }
    & \\ \\

    \theutterance \stepcounter{utterance}  
    & & & \multicolumn{4}{p{0.6\linewidth}}{
        \cellcolor[rgb]{0.9,0.9,0.9}{
            \makecell[{{p{\linewidth}}}]{
                \texttt{\tiny{[GM$\langle$P2]}}
                \texttt{ÜTLE: Oot! Meie eesmärk ei ole viia kõiki objekte ühte lahtrisse (2,2), sest see on füüsiliselt võimatu (üks lahter mahub ainult ühe objekti). Meie eesmärk on, et meie C oleks samas kohas kui sinu C, meie L samas kohas kui sinu L, ja meie P samas kohas kui sinu P.} \\
\\ 
\texttt{Sinu viimases sõnumis mainisid, et sinu objektid on veerus 2: C(2,1), L(2,5), P(2,7).} \\
\texttt{Minu objektide algasukohad olid: C(3,2), L(1,7), P(2,5).} \\
\texttt{Olen juba liigutanud oma C kohale (2,2) ja P kohale (2,4).} \\
\\ 
\texttt{Kuna sinu C on nüüd (2,2) (pärast sinu liikumist), siis minu C on juba õiges kohas (2,2).} \\
\texttt{Sinu L on (2,5) ja sinu P on (2,7).} \\
\\ 
\texttt{Palun kinnita, kas sinu praegused asukohad on tõesti:} \\
\texttt{C: (2,2)} \\
\texttt{L: (2,5)} \\
\texttt{P: (2,7)} \\
\\ 
\texttt{Kui jah, siis peaksin liigutama oma L kohale (2,5) ja oma P kohale (2,7).} \\
\texttt{Kontrollin: Kas (2,5) on tühi? Praegu on seal mu P, mis on juba liigutatud (2,4){-}sse. Seega (2,5) peaks olema tühi.} \\
\texttt{Kontrollin: Kas (2,7) on tühi? Seal on mu L. Seega (2,7) ei ole tühi, kui L on veel seal.} \\
\\ 
\texttt{Olen segaduses. Palun selgita, millised on sinu täpsed praegused asukohad ja millised on minu õiged sihtkohad vastavalt sinu asukohtadele. Me peame lihtsalt kattuma.} \\
\\ 
\texttt{Minu C on (2,2). Sinu C on (2,2). OK.} \\
\texttt{Minu P on (2,4). Sinu P on (2,7). Ei kattu.} \\
\texttt{Minu L on (1,7). Sinu L on (2,5). Ei kattu.} \\
\\ 
\texttt{Ma liigutan oma P kohale (2,7) ja oma L kohale (2,5).} \\
\texttt{Aga oota, (2,7) on praegu tühi? Jah, L on veel (1,7).} \\
\texttt{(2,5) on praegu tühi? Jah, P on juba (2,4).} \\
\\ 
\texttt{Seega järgmine käik: LIIGUTA: L, (2, 5)} \\
            }
        }
    }
     \\ \\

    \theutterance \stepcounter{utterance}  
    & & & \multicolumn{2}{p{0.3\linewidth}}{
        \cellcolor[rgb]{0.9,0.9,0.9}{
            \makecell[{{p{\linewidth}}}]{
                \texttt{\tiny{[GM$|$GM]}}
                \texttt{ÜTLE: Oot! Meie eesmärk ei ole viia kõiki objekte ühte lahtrisse (2,2), sest see on füüsiliselt võimatu (üks lahter mahub ainult ühe objekti). Meie eesmärk on, et meie C oleks samas kohas kui sinu C, meie L samas kohas kui sinu L, ja meie P samas kohas kui sinu P.} \\
\\ 
\texttt{Sinu viimases sõnumis mainisid, et sinu objektid on veerus 2: C(2,1), L(2,5), P(2,7).} \\
\texttt{Minu objektide algasukohad olid: C(3,2), L(1,7), P(2,5).} \\
\texttt{Olen juba liigutanud oma C kohale (2,2) ja P kohale (2,4).} \\
\\ 
\texttt{Kuna sinu C on nüüd (2,2) (pärast sinu liikumist), siis minu C on juba õiges kohas (2,2).} \\
\texttt{Sinu L on (2,5) ja sinu P on (2,7).} \\
\\ 
\texttt{Palun kinnita, kas sinu praegused asukohad on tõesti:} \\
\texttt{C: (2,2)} \\
\texttt{L: (2,5)} \\
\texttt{P: (2,7)} \\
\\ 
\texttt{Kui jah, siis peaksin liigutama oma L kohale (2,5) ja oma P kohale (2,7).} \\
\texttt{Kontrollin: Kas (2,5) on tühi? Praegu on seal mu P, mis on juba liigutatud (2,4){-}sse. Seega (2,5) peaks olema tühi.} \\
\texttt{Kontrollin: Kas (2,7) on tühi? Seal on mu L. Seega (2,7) ei ole tühi, kui L on veel seal.} \\
\\ 
\texttt{Olen segaduses. Palun selgita, millised on sinu täpsed praegused asukohad ja millised on minu õiged sihtkohad vastavalt sinu asukohtadele. Me peame lihtsalt kattuma.} \\
\\ 
\texttt{Minu C on (2,2). Sinu C on (2,2). OK.} \\
\texttt{Minu P on (2,4). Sinu P on (2,7). Ei kattu.} \\
\texttt{Minu L on (1,7). Sinu L on (2,5). Ei kattu.} \\
\\ 
\texttt{Ma liigutan oma P kohale (2,7) ja oma L kohale (2,5).} \\
\texttt{Aga oota, (2,7) on praegu tühi? Jah, L on veel (1,7).} \\
\texttt{(2,5) on praegu tühi? Jah, P on juba (2,4).} \\
\\ 
\texttt{Seega järgmine käik: LIIGUTA: L, (2, 5)} \\
            }
        }
    }
    & & \\ \\

    \theutterance \stepcounter{utterance}  
    & & & \multicolumn{2}{p{0.3\linewidth}}{
        \cellcolor[rgb]{0.9,0.9,0.9}{
            \makecell[{{p{\linewidth}}}]{
                \texttt{\tiny{[GM$|$GM]}}
                \texttt{Invalid response format: ÜTLE: Oot! Meie eesmärk ei ole viia kõiki objekte ühte lahtrisse (2,2), sest see on füüsiliselt võimatu (üks lahter mahub ainult ühe objekti). Meie eesmärk on, et meie C oleks samas kohas kui sinu C, meie L samas kohas kui sinu L, ja meie P samas kohas kui sinu P.} \\
\\ 
\texttt{Sinu viimases sõnumis mainisid, et sinu objektid on veerus 2: C(2,1), L(2,5), P(2,7).} \\
\texttt{Minu objektide algasukohad olid: C(3,2), L(1,7), P(2,5).} \\
\texttt{Olen juba liigutanud oma C kohale (2,2) ja P kohale (2,4).} \\
\\ 
\texttt{Kuna sinu C on nüüd (2,2) (pärast sinu liikumist), siis minu C on juba õiges kohas (2,2).} \\
\texttt{Sinu L on (2,5) ja sinu P on (2,7).} \\
\\ 
\texttt{Palun kinnita, kas sinu praegused asukohad on tõesti:} \\
\texttt{C: (2,2)} \\
\texttt{L: (2,5)} \\
\texttt{P: (2,7)} \\
\\ 
\texttt{Kui jah, siis peaksin liigutama oma L kohale (2,5) ja oma P kohale (2,7).} \\
\texttt{Kontrollin: Kas (2,5) on tühi? Praegu on seal mu P, mis on juba liigutatud (2,4){-}sse. Seega (2,5) peaks olema tühi.} \\
\texttt{Kontrollin: Kas (2,7) on tühi? Seal on mu L. Seega (2,7) ei ole tühi, kui L on veel seal.} \\
\\ 
\texttt{Olen segaduses. Palun selgita, millised on sinu täpsed praegused asukohad ja millised on minu õiged sihtkohad vastavalt sinu asukohtadele. Me peame lihtsalt kattuma.} \\
\\ 
\texttt{Minu C on (2,2). Sinu C on (2,2). OK.} \\
\texttt{Minu P on (2,4). Sinu P on (2,7). Ei kattu.} \\
\texttt{Minu L on (1,7). Sinu L on (2,5). Ei kattu.} \\
\\ 
\texttt{Ma liigutan oma P kohale (2,7) ja oma L kohale (2,5).} \\
\texttt{Aga oota, (2,7) on praegu tühi? Jah, L on veel (1,7).} \\
\texttt{(2,5) on praegu tühi? Jah, P on juba (2,4).} \\
\\ 
\texttt{Seega järgmine käik: LIIGUTA: L, (2, 5)} \\
            }
        }
    }
    & & \\ \\

    \theutterance \stepcounter{utterance}  
    & & \multicolumn{4}{p{0.6\linewidth}}{
        \cellcolor[rgb]{0.9,0.9,0.9}{
            \makecell[{{p{\linewidth}}}]{
                \texttt{\tiny{[GM$\rangle$P2]}}
                \texttt{Karistus: Teie sõnum sisaldab rohkem kui üht käsku!} \\
\\ 
\texttt{Veendu, et sinu vastus sisaldab ainult kas `ÜTLE: <SÕNUM>` või `LIIGUTA: <OBJEKT>, (<X>, <Y>)` ja mitte midagi muud!} \\
\\ 
\texttt{Teil on kollektiivselt kogunenud 8 karistust 8 karistusest.} \\
\texttt{Palun proovi uuesti!} \\
            }
        }
    }
    & \\ \\

    \theutterance \stepcounter{utterance}  
    & & & \multicolumn{4}{p{0.6\linewidth}}{
        \cellcolor[rgb]{0.9,0.9,0.9}{
            \makecell[{{p{\linewidth}}}]{
                \texttt{\tiny{[GM$\langle$P2]}}
                \texttt{ÜTLE: Vabandust, eelmine sõnum oli valesti vormistatud. Meie eesmärk on kattuda: minu C peab olema sinu C juures, minu L sinu L juures ja minu P sinu P juures. Sinu asukohad on C(2,2), L(2,5), P(2,7). Minu C on juba (2,2). Minu P on (2,4) ja minu L on (1,7). Liigutan nüüd oma L kohale (2,5), kuna see on tühi.} \\
            }
        }
    }
     \\ \\

    \theutterance \stepcounter{utterance}  
    & & & \multicolumn{2}{p{0.3\linewidth}}{
        \cellcolor[rgb]{0.9,0.9,0.9}{
            \makecell[{{p{\linewidth}}}]{
                \texttt{\tiny{[GM$|$GM]}}
                \texttt{ÜTLE: Vabandust, eelmine sõnum oli valesti vormistatud. Meie eesmärk on kattuda: minu C peab olema sinu C juures, minu L sinu L juures ja minu P sinu P juures. Sinu asukohad on C(2,2), L(2,5), P(2,7). Minu C on juba (2,2). Minu P on (2,4) ja minu L on (1,7). Liigutan nüüd oma L kohale (2,5), kuna see on tühi.} \\
            }
        }
    }
    & & \\ \\

    \theutterance \stepcounter{utterance}  
    & & \multicolumn{4}{p{0.6\linewidth}}{
        \cellcolor[rgb]{0.9,0.9,0.9}{
            \makecell[{{p{\linewidth}}}]{
                \texttt{\tiny{[P1$\langle$GM]}}
                \texttt{Sinu sõnum on teisele mängijale edastatud.} \\
\\ 
\\ 
\texttt{Sa mängid praegu vooru 6 maksimaalsest 12 voorust.} \\
\texttt{Teil on kollektiivselt kogunenud 8 karistust 8 karistusest.} \\
\texttt{Teine mängija saatis selle sõnumi:} \\
\\ 
\\ 
\texttt{"Vabandust, eelmine sõnum oli valesti vormistatud. Meie eesmärk on kattuda: minu C peab olema sinu C juures, minu L sinu L juures ja minu P sinu P juures. Sinu asukohad on C(2,2), L(2,5), P(2,7). Minu C on juba (2,2). Minu P on (2,4) ja minu L on (1,7). Liigutan nüüd oma L kohale (2,5), kuna see on tühi."} \\
\\ 
\\ 
\texttt{Mis on sinu järgmine käsk?} \\
            }
        }
    }
    & \\ \\

    \theutterance \stepcounter{utterance}  
    & \multicolumn{4}{p{0.6\linewidth}}{
        \cellcolor[rgb]{0.9,0.9,0.9}{
            \makecell[{{p{\linewidth}}}]{
                \texttt{\tiny{[P1$\rangle$GM]}}
                \texttt{LIIGUTA: C, (2, 2)} \\
            }
        }
    }
    & & \\ \\

    \theutterance \stepcounter{utterance}  
    & & & \multicolumn{2}{p{0.3\linewidth}}{
        \cellcolor[rgb]{0.9,0.9,0.9}{
            \makecell[{{p{\linewidth}}}]{
                \texttt{\tiny{[GM$|$GM]}}
                \texttt{LIIGUTA: C, (2, 2)} \\
            }
        }
    }
    & & \\ \\

    \theutterance \stepcounter{utterance}  
    & & & \multicolumn{2}{p{0.3\linewidth}}{
        \cellcolor[rgb]{0.9,0.9,0.9}{
            \makecell[{{p{\linewidth}}}]{
                \texttt{\tiny{[GM$|$GM]}}
                \texttt{C moved to (2, 2)} \\
            }
        }
    }
    & & \\ \\

    \theutterance \stepcounter{utterance}  
    & & \multicolumn{4}{p{0.6\linewidth}}{
        \cellcolor[rgb]{0.9,0.9,0.9}{
            \makecell[{{p{\linewidth}}}]{
                \texttt{\tiny{[GM$\rangle$P2]}}
                \texttt{Sinu sõnum on teisele mängijale edastatud.} \\
\\ 
\\ 
\texttt{Sa mängid praegu vooru 7 maksimaalsest 12 voorust.} \\
\texttt{Teil on kollektiivselt kogunenud 8 karistust 8 karistusest.} \\
\texttt{Teine mängija liigutas objekti oma ruudustikul.} \\
\\ 
\\ 
\texttt{Mis on sinu järgmine käsk?} \\
            }
        }
    }
    & \\ \\

    \theutterance \stepcounter{utterance}  
    & & & \multicolumn{4}{p{0.6\linewidth}}{
        \cellcolor[rgb]{0.9,0.9,0.9}{
            \makecell[{{p{\linewidth}}}]{
                \texttt{\tiny{[GM$\langle$P2]}}
                \texttt{LIIGUTA: L, (2, 5)} \\
            }
        }
    }
     \\ \\

    \theutterance \stepcounter{utterance}  
    & & & \multicolumn{2}{p{0.3\linewidth}}{
        \cellcolor[rgb]{0.9,0.9,0.9}{
            \makecell[{{p{\linewidth}}}]{
                \texttt{\tiny{[GM$|$GM]}}
                \texttt{LIIGUTA: L, (2, 5)} \\
            }
        }
    }
    & & \\ \\

    \theutterance \stepcounter{utterance}  
    & & & \multicolumn{2}{p{0.3\linewidth}}{
        \cellcolor[rgb]{0.9,0.9,0.9}{
            \makecell[{{p{\linewidth}}}]{
                \texttt{\tiny{[GM$|$GM]}}
                \texttt{L moved to (2, 5)} \\
            }
        }
    }
    & & \\ \\

    \theutterance \stepcounter{utterance}  
    & & \multicolumn{4}{p{0.6\linewidth}}{
        \cellcolor[rgb]{0.9,0.9,0.9}{
            \makecell[{{p{\linewidth}}}]{
                \texttt{\tiny{[P1$\langle$GM]}}
                \texttt{Edukalt teisaldatud 'C' asukohta (2,2). Teie uuendatud ruudustik näeb välja selline:} \\
\texttt{```} \\
\texttt{ 1234567} \\
\texttt{╔═══════╗} \\
\texttt{║◌◌◌◌◌◌◌║ 1} \\
\texttt{║◌C◌◌◌◌◌║ 2} \\
\texttt{║◌◌◌◌◌┌─╢ 3} \\
\texttt{║◌◌◌┌─┤◌║ 4} \\
\texttt{║◌L◌│◌├─╢ 5} \\
\texttt{╟───┼─┘◌║ 6} \\
\texttt{║◌P◌│◌◌◌║ 7} \\
\texttt{╚═══╧═══╝} \\
\texttt{```} \\
\texttt{Sa mängid praegu vooru 7 maksimaalsest 12 voorust.} \\
\texttt{Teil on kollektiivselt kogunenud 8 karistust 8 karistusest.} \\
\texttt{Teine mängija liigutas objekti oma ruudustikul.} \\
\\ 
\\ 
\texttt{Mis on sinu järgmine käsk?} \\
            }
        }
    }
    & \\ \\

    \theutterance \stepcounter{utterance}  
    & \multicolumn{4}{p{0.6\linewidth}}{
        \cellcolor[rgb]{0.9,0.9,0.9}{
            \makecell[{{p{\linewidth}}}]{
                \texttt{\tiny{[P1$\rangle$GM]}}
                \texttt{ÜTLE: valmis?} \\
            }
        }
    }
    & & \\ \\

    \theutterance \stepcounter{utterance}  
    & & & \multicolumn{2}{p{0.3\linewidth}}{
        \cellcolor[rgb]{0.9,0.9,0.9}{
            \makecell[{{p{\linewidth}}}]{
                \texttt{\tiny{[GM$|$GM]}}
                \texttt{ÜTLE: valmis?} \\
            }
        }
    }
    & & \\ \\

    \theutterance \stepcounter{utterance}  
    & & \multicolumn{4}{p{0.6\linewidth}}{
        \cellcolor[rgb]{0.9,0.9,0.9}{
            \makecell[{{p{\linewidth}}}]{
                \texttt{\tiny{[GM$\rangle$P2]}}
                \texttt{Edukalt teisaldatud 'L' asukohta (2,5). Teie uuendatud ruudustik näeb välja selline:} \\
\texttt{```} \\
\texttt{ 1234567} \\
\texttt{╔═══════╗} \\
\texttt{║◌◌◌◌◌◌◌║ 1} \\
\texttt{║◌C◌◌◌◌◌║ 2} \\
\texttt{║◌◌◌◌◌┌─╢ 3} \\
\texttt{║◌P◌┌─┤◌║ 4} \\
\texttt{║◌L◌│◌├─╢ 5} \\
\texttt{╟───┼─┘◌║ 6} \\
\texttt{║◌◌◌│◌◌◌║ 7} \\
\texttt{╚═══╧═══╝} \\
\texttt{```} \\
\texttt{Sa mängid praegu vooru 8 maksimaalsest 12 voorust.} \\
\texttt{Teil on kollektiivselt kogunenud 8 karistust 8 karistusest.} \\
\texttt{Teine mängija saatis selle sõnumi:} \\
\\ 
\\ 
\texttt{"valmis?"} \\
\\ 
\\ 
\texttt{Mis on sinu järgmine käsk?} \\
            }
        }
    }
    & \\ \\

    \theutterance \stepcounter{utterance}  
    & & & \multicolumn{4}{p{0.6\linewidth}}{
        \cellcolor[rgb]{0.9,0.9,0.9}{
            \makecell[{{p{\linewidth}}}]{
                \texttt{\tiny{[GM$\langle$P2]}}
                \texttt{ÜTLE: Ei ole veel valmis. Mul on vaja veel liigutada oma P objekt kohale (2,7), et see kattuks sinu P{-}ga. Palun oota.} \\
            }
        }
    }
     \\ \\

    \theutterance \stepcounter{utterance}  
    & & & \multicolumn{2}{p{0.3\linewidth}}{
        \cellcolor[rgb]{0.9,0.9,0.9}{
            \makecell[{{p{\linewidth}}}]{
                \texttt{\tiny{[GM$|$GM]}}
                \texttt{ÜTLE: Ei ole veel valmis. Mul on vaja veel liigutada oma P objekt kohale (2,7), et see kattuks sinu P{-}ga. Palun oota.} \\
            }
        }
    }
    & & \\ \\

    \theutterance \stepcounter{utterance}  
    & & \multicolumn{4}{p{0.6\linewidth}}{
        \cellcolor[rgb]{0.9,0.9,0.9}{
            \makecell[{{p{\linewidth}}}]{
                \texttt{\tiny{[P1$\langle$GM]}}
                \texttt{Sinu sõnum on teisele mängijale edastatud.} \\
\\ 
\\ 
\texttt{Sa mängid praegu vooru 8 maksimaalsest 12 voorust.} \\
\texttt{Teil on kollektiivselt kogunenud 8 karistust 8 karistusest.} \\
\texttt{Teine mängija saatis selle sõnumi:} \\
\\ 
\\ 
\texttt{"Ei ole veel valmis. Mul on vaja veel liigutada oma P objekt kohale (2,7), et see kattuks sinu P{-}ga. Palun oota."} \\
\\ 
\\ 
\texttt{Mis on sinu järgmine käsk?} \\
            }
        }
    }
    & \\ \\

    \theutterance \stepcounter{utterance}  
    & \multicolumn{4}{p{0.6\linewidth}}{
        \cellcolor[rgb]{0.9,0.9,0.9}{
            \makecell[{{p{\linewidth}}}]{
                \texttt{\tiny{[P1$\rangle$GM]}}
                \texttt{ÜTLE: Okei, ootan. Palun anna teada, kui oled valmis.} \\
            }
        }
    }
    & & \\ \\

    \theutterance \stepcounter{utterance}  
    & & & \multicolumn{2}{p{0.3\linewidth}}{
        \cellcolor[rgb]{0.9,0.9,0.9}{
            \makecell[{{p{\linewidth}}}]{
                \texttt{\tiny{[GM$|$GM]}}
                \texttt{ÜTLE: Okei, ootan. Palun anna teada, kui oled valmis.} \\
            }
        }
    }
    & & \\ \\

    \theutterance \stepcounter{utterance}  
    & & \multicolumn{4}{p{0.6\linewidth}}{
        \cellcolor[rgb]{0.9,0.9,0.9}{
            \makecell[{{p{\linewidth}}}]{
                \texttt{\tiny{[GM$\rangle$P2]}}
                \texttt{Sinu sõnum on teisele mängijale edastatud.} \\
\\ 
\\ 
\texttt{Sa mängid praegu vooru 9 maksimaalsest 12 voorust.} \\
\texttt{Teil on kollektiivselt kogunenud 8 karistust 8 karistusest.} \\
\texttt{Teine mängija saatis selle sõnumi:} \\
\\ 
\\ 
\texttt{"Okei, ootan. Palun anna teada, kui oled valmis."} \\
\\ 
\\ 
\texttt{Mis on sinu järgmine käsk?} \\
            }
        }
    }
    & \\ \\

    \theutterance \stepcounter{utterance}  
    & & & \multicolumn{4}{p{0.6\linewidth}}{
        \cellcolor[rgb]{0.9,0.9,0.9}{
            \makecell[{{p{\linewidth}}}]{
                \texttt{\tiny{[GM$\langle$P2]}}
                \texttt{LIIGUTA: P, (2, 7)} \\
            }
        }
    }
     \\ \\

    \theutterance \stepcounter{utterance}  
    & & & \multicolumn{2}{p{0.3\linewidth}}{
        \cellcolor[rgb]{0.9,0.9,0.9}{
            \makecell[{{p{\linewidth}}}]{
                \texttt{\tiny{[GM$|$GM]}}
                \texttt{LIIGUTA: P, (2, 7)} \\
            }
        }
    }
    & & \\ \\

    \theutterance \stepcounter{utterance}  
    & & & \multicolumn{2}{p{0.3\linewidth}}{
        \cellcolor[rgb]{0.9,0.9,0.9}{
            \makecell[{{p{\linewidth}}}]{
                \texttt{\tiny{[GM$|$GM]}}
                \texttt{P moved to (2, 7)} \\
            }
        }
    }
    & & \\ \\

    \theutterance \stepcounter{utterance}  
    & & \multicolumn{4}{p{0.6\linewidth}}{
        \cellcolor[rgb]{0.9,0.9,0.9}{
            \makecell[{{p{\linewidth}}}]{
                \texttt{\tiny{[P1$\langle$GM]}}
                \texttt{Sinu sõnum on teisele mängijale edastatud.} \\
\\ 
\\ 
\texttt{Sa mängid praegu vooru 9 maksimaalsest 12 voorust.} \\
\texttt{Teil on kollektiivselt kogunenud 8 karistust 8 karistusest.} \\
\texttt{Teine mängija liigutas objekti oma ruudustikul.} \\
\\ 
\\ 
\texttt{Mis on sinu järgmine käsk?} \\
            }
        }
    }
    & \\ \\

    \theutterance \stepcounter{utterance}  
    & \multicolumn{4}{p{0.6\linewidth}}{
        \cellcolor[rgb]{0.9,0.9,0.9}{
            \makecell[{{p{\linewidth}}}]{
                \texttt{\tiny{[P1$\rangle$GM]}}
                \texttt{ÜTLE: valmis?} \\
            }
        }
    }
    & & \\ \\

    \theutterance \stepcounter{utterance}  
    & & & \multicolumn{2}{p{0.3\linewidth}}{
        \cellcolor[rgb]{0.9,0.9,0.9}{
            \makecell[{{p{\linewidth}}}]{
                \texttt{\tiny{[GM$|$GM]}}
                \texttt{ÜTLE: valmis?} \\
            }
        }
    }
    & & \\ \\

    \theutterance \stepcounter{utterance}  
    & & \multicolumn{4}{p{0.6\linewidth}}{
        \cellcolor[rgb]{0.9,0.9,0.9}{
            \makecell[{{p{\linewidth}}}]{
                \texttt{\tiny{[GM$\rangle$P2]}}
                \texttt{Edukalt teisaldatud 'P' asukohta (2,7). Teie uuendatud ruudustik näeb välja selline:} \\
\texttt{```} \\
\texttt{ 1234567} \\
\texttt{╔═══════╗} \\
\texttt{║◌◌◌◌◌◌◌║ 1} \\
\texttt{║◌C◌◌◌◌◌║ 2} \\
\texttt{║◌◌◌◌◌┌─╢ 3} \\
\texttt{║◌◌◌┌─┤◌║ 4} \\
\texttt{║◌L◌│◌├─╢ 5} \\
\texttt{╟───┼─┘◌║ 6} \\
\texttt{║◌P◌│◌◌◌║ 7} \\
\texttt{╚═══╧═══╝} \\
\texttt{```} \\
\texttt{Sa mängid praegu vooru 10 maksimaalsest 12 voorust.} \\
\texttt{Teil on kollektiivselt kogunenud 8 karistust 8 karistusest.} \\
\texttt{Teine mängija saatis selle sõnumi:} \\
\\ 
\\ 
\texttt{"valmis?"} \\
\\ 
\\ 
\texttt{Mis on sinu järgmine käsk?} \\
            }
        }
    }
    & \\ \\

    \theutterance \stepcounter{utterance}  
    & & & \multicolumn{4}{p{0.6\linewidth}}{
        \cellcolor[rgb]{0.9,0.9,0.9}{
            \makecell[{{p{\linewidth}}}]{
                \texttt{\tiny{[GM$\langle$P2]}}
                \texttt{ÜTLE: valmis!} \\
            }
        }
    }
     \\ \\

    \theutterance \stepcounter{utterance}  
    & & & \multicolumn{2}{p{0.3\linewidth}}{
        \cellcolor[rgb]{0.9,0.9,0.9}{
            \makecell[{{p{\linewidth}}}]{
                \texttt{\tiny{[GM$|$GM]}}
                \texttt{ÜTLE: valmis!} \\
            }
        }
    }
    & & \\ \\

    \theutterance \stepcounter{utterance}  
    & & & \multicolumn{2}{p{0.3\linewidth}}{
        \cellcolor[rgb]{0.9,0.9,0.9}{
            \makecell[{{p{\linewidth}}}]{
                \texttt{\tiny{[GM$|$GM]}}
                \texttt{true} \\
            }
        }
    }
    & & \\ \\

    \theutterance \stepcounter{utterance}  
    & & & \multicolumn{2}{p{0.3\linewidth}}{
        \cellcolor[rgb]{0.9,0.9,0.9}{
            \makecell[{{p{\linewidth}}}]{
                \texttt{\tiny{[GM$|$GM]}}
                \texttt{Initial states:} \\
\texttt{```} \\
\texttt{Player 1:} \\
\texttt{ 1234567} \\
\texttt{╔═══════╗} \\
\texttt{║◌◌◌◌◌◌◌║ 1} \\
\texttt{║◌C◌◌◌◌◌║ 2} \\
\texttt{║◌◌◌◌◌┌─╢ 3} \\
\texttt{║◌◌◌┌─┤◌║ 4} \\
\texttt{║◌◌◌│L├─╢ 5} \\
\texttt{╟───┼─┘◌║ 6} \\
\texttt{║◌◌◌│◌◌P║ 7} \\
\texttt{╚═══╧═══╝} \\
\\ 
\\ 
\texttt{Player 2:} \\
\texttt{ 1234567} \\
\texttt{╔═══════╗} \\
\texttt{║◌◌◌◌◌◌◌║ 1} \\
\texttt{║◌◌C◌◌◌◌║ 2} \\
\texttt{║◌◌◌◌◌┌─╢ 3} \\
\texttt{║◌◌◌┌─┤◌║ 4} \\
\texttt{║◌P◌│◌├─╢ 5} \\
\texttt{╟───┼─┘◌║ 6} \\
\texttt{║L◌◌│◌◌◌║ 7} \\
\texttt{╚═══╧═══╝} \\
\\ 
\texttt{```} \\
            }
        }
    }
    & & \\ \\

    \theutterance \stepcounter{utterance}  
    & & & \multicolumn{2}{p{0.3\linewidth}}{
        \cellcolor[rgb]{0.9,0.9,0.9}{
            \makecell[{{p{\linewidth}}}]{
                \texttt{\tiny{[GM$|$GM]}}
                \texttt{End states:} \\
\texttt{```} \\
\texttt{Player 1:} \\
\texttt{ 1234567} \\
\texttt{╔═══════╗} \\
\texttt{║◌◌◌◌◌◌◌║ 1} \\
\texttt{║◌C◌◌◌◌◌║ 2} \\
\texttt{║◌◌◌◌◌┌─╢ 3} \\
\texttt{║◌◌◌┌─┤◌║ 4} \\
\texttt{║◌L◌│◌├─╢ 5} \\
\texttt{╟───┼─┘◌║ 6} \\
\texttt{║◌P◌│◌◌◌║ 7} \\
\texttt{╚═══╧═══╝} \\
\\ 
\\ 
\texttt{Player 2:} \\
\texttt{ 1234567} \\
\texttt{╔═══════╗} \\
\texttt{║◌◌◌◌◌◌◌║ 1} \\
\texttt{║◌C◌◌◌◌◌║ 2} \\
\texttt{║◌◌◌◌◌┌─╢ 3} \\
\texttt{║◌◌◌┌─┤◌║ 4} \\
\texttt{║◌L◌│◌├─╢ 5} \\
\texttt{╟───┼─┘◌║ 6} \\
\texttt{║◌P◌│◌◌◌║ 7} \\
\texttt{╚═══╧═══╝} \\
\\ 
\texttt{```} \\
            }
        }
    }
    & & \\ \\

    \theutterance \stepcounter{utterance}  
    & & & \multicolumn{2}{p{0.3\linewidth}}{
        \cellcolor[rgb]{0.9,0.9,0.9}{
            \makecell[{{p{\linewidth}}}]{
                \texttt{\tiny{[GM$|$GM]}}
                \texttt{* success: True} \\
\texttt{* lose: False} \\
\texttt{* aborted: False} \\
\texttt{{-}{-}{-}{-}{-}{-}{-}} \\
\texttt{* Shifts: 6.00} \\
\texttt{* Max Shifts: 4.00} \\
\texttt{* Min Shifts: 2.00} \\
\texttt{* End Distance Sum: 0.00} \\
\texttt{* Init Distance Sum: 10.86} \\
\texttt{* Expected Distance Sum: 12.57} \\
\texttt{* Penalties: 8.00} \\
\texttt{* Max Penalties: 8.00} \\
\texttt{* Rounds: 10.00} \\
\texttt{* Max Rounds: 12.00} \\
\texttt{* Object Count: 3.00} \\
            }
        }
    }
    & & \\ \\

\end{supertabular}
}

\end{document}
